% === IV. SIMULATION ARCHITECTURE AND DATA PIPELINE ===========================
\section{Simulation Architecture and Data Pipeline}
\label{sec:platform}

The identification pipeline is exercised end to end inside a single simulation
architecture: a Formula Student vehicle model in CARLA acting as the plant, a
C\texttt{++} ROS~2 communication layer that drives it and republishes its state, and an
identification stack that consumes that state. The architecture is described
here because two of its properties---per-wheel ground truth and
simulation-time stamping---are what make the identifiability claims of
Section~\ref{sec:method} checkable rather than merely arguable.

\subsection{Choice of simulator}
\label{subsec:whycarla}

CARLA~\cite{carla2017} was selected against four requirements that the
identification problem of Section~\ref{sec:problem} imposes on a simulator.

\emph{Per-wheel ground truth.} The claim to be checked is that the recovered
Magic Formula reproduces the lateral force the plant actually produced, so the
simulator must expose per-wheel slip angle, normal load and tire force---not
only chassis states. CARLA's PhysX wheeled-vehicle solver computes these
quantities internally and its client API exports them, which is what makes the
tire forces feedback channel of Section~\ref{subsec:bridge} possible. 
Cone-oriented Formula Student simulators.

\emph{Deterministic, simulation-time-stamped execution.} Persistent-excitation
and residual-learning arguments are only checkable if every sample carries a
consistent clock. CARLA's synchronous mode with a fixed physics step
(\SI{0.033}{\second} here) decouples data timing from wall-clock load; all
identification data is stamped from the simulation clock
(Appendix, Section~A7).

\emph{Full-scale vehicle and racing-track support.}
The target class is a full-scale Formula Student car on a closed circuit.
CARLA accepts custom vehicle assets here the Formula AI vehicle on the
silverstone map and offers a client that the bridge links against directly, 
avoiding a Python round trip in the control loop. Surveys of open self-driving simulators place
CARLA (with the now-discontinued LGSVL) at the state of the art for end-to-end
stack testing on exactly these grounds~\cite{kaur2021simsurvey}.

\emph{A plant whose tire behaviour is \emph{not} a Magic Formula.} CARLA's
native dynamics are a known fidelity gap: R-CARLA~\cite{rcarla2025} replaces
the dynamics model to shrink a measured sim-to-real gap, and Sagmeister
\emph{et al.}~\cite{simfidelity2026} quantify where model simplification
changes control conclusions. For this study that gap is a feature rather than a
defect. The plant's lateral force comes from PhysX's default tire
model~\cite{physxvehicle} and not from \eqref{eq:mf}, so the simulator holds no
Pacejka coefficients for the identifier to trivially recover, and the two model
families differ in shape and not only in parameters
(Section~\ref{subsec:planttire}). That is the same position the method occupies
on a real vehicle, with one advantage a real vehicle does not offer: the
plant's curve is available in closed form from the simulator's own
configuration, so the identification is scored against an exact target rather
than against a fit. The limits of this choice are revisited in
Section~\ref{sec:disc}.

\subsection{Plant: Formula AI vehicle in CARLA}
\label{subsec:platform}

\begin{figure}[t]
\centering
\includegraphics[width=0.86\columnwidth]{fig2.png}
\caption{The Formula AI vehicle model, on the start point of the
silverstone map in CARLA. Its lateral forces are produced by the engine's
PhysX wheeled-vehicle tire model~\cite{physxvehicle} rather than by a Magic
Formula, which is the curve of Table~\ref{tab:ref} the identifier is scored
against. Measured geometry and mass are given in Table~\ref{tab:vehicle}.}
\label{fig:vehicle}
\end{figure}

The plant is the driverless Formula Student--class vehicle model, instantiated in
CARLA~\cite{carla2017} (server 0.9.16, Unreal Engine 4.26) on the
\texttt{silverstone} map: a four-wheeled, front-steered, all-wheel-drive
rigid-body vehicle whose lateral forces are generated by the PhysX
wheeled-vehicle tire model~\cite{physxvehicle} under a configured friction
coefficient and lateral stiffness, not by a Magic Formula. The simulator
therefore has no ``true'' Pacejka coefficients to recover; what it does have is
its own force law in closed form, and Section~\ref{subsec:planttire} states it
and checks it against the per-wheel telemetry, so the identification of
Section~\ref{sec:results} is scored against the plant's curve rather than
against a fit to it.

\begin{table}[t]
\caption{Vehicle configuration and measured plant properties. Configured values
are set in the bridge's \texttt{carla\_interface\_\allowbreak config.yaml};
measured values are obtained from the simulator's per-wheel telemetry.}
\label{tab:vehicle}
\centering
\footnotesize
\setlength{\tabcolsep}{4pt}
\begin{tabular}{@{}llll@{}}
\toprule
Quantity  & Value & Unit\\
\midrule
Body mass                      & 240.0 & \si{\kilo\gram}\\
Static wheel load (total)      & 2644.5 & \si{\newton}\\
Unsprung mass                  & 269.8 & \si{\kilo\gram}\\
Front/rear static split        & 51.2 / 48.8 & \si{\percent}\\
Wheelbase $L$                  & 1.5334 & \si{\meter}\\
$l_{f}$ / $l_{r}$              & 0.7381 / 0.7953 & \si{\meter}\\
Yaw inertia $I_{z}$\textsuperscript{b} & 51.1 & \si{\kilo\gram\meter\squared}\\
Wheel radius                   & 0.25 & \si{\meter}\\
Max road-wheel angle           & 16.0 & \si{\degree}\\
Wheel friction                 & 1.05 & ---\\
Realised peak axle friction    & 1.00--1.05 & ---\\
Drag coefficient               & 0.426 & ---\\
\bottomrule
\end{tabular}

\vspace{2pt}
{\raggedright\footnotesize\textsuperscript{a}Total static wheel load divided by
$g$; see text. \textsuperscript{b}Assumed; see
Section~\ref{subsec:assumptions}. \textsuperscript{c}Configured wheel friction
multiplied by the road surface's own coefficient, as reported per wheel by
\texttt{feedback/\allowbreak tire\_forces}.\par}
\end{table}

Three groups of rows in Table~\ref{tab:vehicle} are load-bearing for the
identification, and all three were wrong in the inherited configuration.

First, mass. The total static normal load the wheels report is
\SI{2644.5}{\newton}, i.e.\ \SI{269.8}{\kilo\gram} against the
\SI{240.0}{\kilo\gram} configured as the body mass---\SI{12.4}{\percent} more.
The difference is not noise: \SI{24.3}{\kilo\gram} of it is unsprung wheel mass
that PhysX carries outside \texttt{vehicle.physics.mass}, measured by sweeping
the configured mass over \num{240}, \num{400} and \SI{600}{\kilo\gram} and observing a
constant offset, and the remainder is standing suspension compression. Since
\eqref{eq:ssforces} and \eqref{eq:newtoneuler} are both linear in $m$, using
the configured value biases every recovered force---and therefore the
identified $D$ and the warm-started $\mu$---low by that same
\SI{12}{\percent}. The identifier now takes $m$ from the measured static load.

Second, friction. The wheel friction coefficient set in the configuration
(\num{1.5}) is not the coefficient the physics step uses: the server multiplies
it by the road surface's own coefficient, a constant factor of \num{0.70} on
this map, and the product (\num{1.05}) is what appears per wheel in the
telemetry's \texttt{tire\_friction} field and what bounds the achievable axle
force. It matches the realised peak axle friction of \numrange{1.00}{1.05}. An
identifier---or a friction-ellipse check---tuned against the configuration file
converges on a target \SI{43}{\percent} too high.

Third, steering. The achieved road-wheel angle is \num{0.76} of the commanded
angle, so slip angles built from the command are systematically too large
(Section~\ref{subsec:measurement}).

\subsection{The plant's tire curve, in closed form}
\label{subsec:planttire}

The identification target is worth stating exactly, because it is knowable
here. PhysX computes the lateral force of each wheel with
\texttt{PxVehicle\-Compute\-Tire\-Force\-Default}~\cite{physxvehicle}; at zero
longitudinal slip and zero camber that expression collapses to
\begin{equation}
F_{y}(\alpha)=\mathrm{sgn}(\alpha)\,\mu F_{z}\,S_{1}(K),\qquad
K=\frac{C_{\mathrm{lat}}|\tan\alpha|}{\mu F_{z}},
\label{eq:physx}
\end{equation}
with $C_{\mathrm{lat}}=F_{z}^{\mathrm{rest}}\,k_{y}\,
S_{1}\!\big(3\bar F_{z}/k_{x}\big)$, $\bar F_{z}=F_{z}/F_{z}^{\mathrm{rest}}$
the normalised load, $(k_{y},k_{x})$ the wheel's lateral-stiffness parameters,
and $S_{1}(K)=\min(1,\,K-K^{2}/3+K^{3}/27)$ the engine's smoothing function.

Two properties of \eqref{eq:physx} shape everything downstream. It saturates at
$\mu F_{z}$, first reached at $K=3$, and it \emph{never falls off} past that
peak---so a Magic Formula, whose $E$ and post-peak decay exist precisely to
describe a falling branch, can match it up to the peak and not beyond. And its
scale is set entirely by $\mu$ and $C_{\mathrm{lat}}$, both of which the
simulator publishes, so the curve of Table~\ref{tab:ref} is read off the plant
rather than fitted to it.

\begin{table}[t]
\caption{The plant's lateral tire curve. Upper block: the parameters of
\eqref{eq:physx}, read from the simulator's own wheel configuration and
telemetry. Lower block: the Magic Formula prior of \eqref{eq:mf} used as the
identifier's starting point and frozen regularisation target, per axle
$[B,C,D,E]$.}
\label{tab:ref}
\centering
\footnotesize
\setlength{\tabcolsep}{4pt}
\begin{tabular}{@{}lll@{}}
\toprule
Quantity & Value & Unit\\
\midrule
Effective friction $\mu$          & 1.05 & ---\\
Lateral stiffness $k_{y}$         & 17.0 & \si{\per\radian}\\
Max normalised load $k_{x}$       & 2.0 & ---\\
Initial slope $\partial F_{y}/\partial\alpha$ & 14.80\,$F_{z}$ & \si{\newton\per\radian}\\
Peak reached at                   & 11.5 & \si{\degree}\\
\midrule
$\Phi_{p,f}$ prior & $[7.076,\,1.346,\,1.009,\,-2.000]$ & ---\\
$\Phi_{p,r}$ prior & $[7.873,\,1.383,\,1.002,\,-1.024]$ & ---\\
\bottomrule
\end{tabular}
\end{table}

That the closed form is the right one was checked rather than assumed. On 1804
ground-truth wheel samples from a live lap ($P_{99}(|\alpha|)=\SI{3.66}{\degree}$),
each predicted from its own measured $\alpha$, $F_{z}$ and $\mu$,
\eqref{eq:physx} returns $R^{2}=\num{0.99}$ with \SI{16.1}{\newton} RMSE on the
front axle and \SI{14.7}{\newton} on the rear, at biases of \SI{+2.0}{\newton}
and \SI{+0.8}{\newton}. Section~\ref{subsec:corrected} uses this curve to
separate the error of the identified fit from the error of representing
\eqref{eq:physx} in the form \eqref{eq:mf} that the controller requires.

\begin{figure*}[t]
\centering
\resizebox{\textwidth}{!}{% =============================================================================
% Fig. 3 -- Carla_ASU_Bridge architecture.
%
% Vector TikZ, drawn at the paper's own font sizes. Included by
% sections/method.tex inside a figure* (full text width).
%
% Reading order, left to right: CARLA server -> bridge lifecycle node ->
% ROS 2 topic layer -> autonomy stack.
%   * orange  = the ground-truth path (feedback/tire_forces). It leaves the
%               simulator and reaches only the scoring/validation nodes; it is
%               deliberately NOT wired into the identification loop.
%   * blue    = mpc/update_params, the gated model interface. Drawn inside the
%               autonomy stack because it is a service between two of this
%               workspace's nodes, not part of the bridge's topic layer.
%
% Nodes are stacked with relative positioning (below=... of ...) so that a
% change of wording cannot make two boxes overlap.
% =============================================================================

\begin{tikzpicture}[
  font=\scriptsize,
  >={Stealth[length=1.7mm,width=1.5mm]},
  % --- block styles ---
  colhead/.style  = {font=\scriptsize\bfseries, align=center},
  blk/.style      = {draw, rounded corners=1pt, align=left, inner sep=2.4pt,
                     text width=#1, fill=white, line width=0.35pt},
  srv/.style      = {blk=26mm, fill=black!5},
  brg/.style      = {blk=33mm, fill=black!3},
  top/.style      = {blk=36mm, font=\scriptsize\ttfamily, fill=white},
  gt/.style       = {top, fill=orange!20, draw=orange!80!black, line width=0.7pt},
  cons/.style     = {blk=31mm, fill=black!3},
  note/.style     = {font=\tiny, align=left, inner sep=1.5pt, text width=36mm},
  % --- edge styles ---
  flow/.style     = {->, line width=0.5pt, black!65},
  gtflow/.style   = {->, line width=0.9pt, orange!80!black},
  gateflow/.style = {->, line width=0.9pt, blue!60!black},
  lbl/.style      = {font=\tiny, inner sep=1.2pt, fill=white, midway}
]

\def\gap{3.2mm}   % vertical gap between stacked boxes

% ============================ 1. CARLA server ============================
\node[srv, anchor=north west] (ego) at (0,0)
  {\textbf{Ego vehicle}\\[1pt]
   {\ttfamily vehicle.vehicle.}\\{\ttfamily asurt\_fsai}\\[1pt]
   PhysX wheeled-vehicle solver,\\per-wheel tire friction};
\node[srv, below=\gap of ego]  (sens)
  {\textbf{Sensor actors}\\[1pt] IMU, GNSS, cameras, LiDAR};
\node[srv, below=\gap of sens] (tick)
  {\textbf{World tick}\\[1pt] synchronous, $\Delta t=0.033$\,s};
\node[srv, below=\gap of tick] (wt)
  {\textbf{Wheel telemetry}\\[1pt]
   per-wheel slip angle,\\normal load, lateral force};

\begin{scope}[on background layer]
  \node[draw, dashed, rounded corners=2pt, inner sep=3.6pt, fill=black!2,
        fit=(ego)(sens)(tick)(wt)] (carla) {};
\end{scope}
\node[colhead, above=1.4mm of carla.north]
  {CARLA server\\{\scriptsize\mdseries UE 4.26 / PhysX}};

% ======================== 2. Carla_ASU_Bridge node ========================
\node[brg, anchor=north west] (ctrl) at (5.75,0)
  {\textbf{Control application}\\[1pt]
   forwards speed and steering to\\CARLA's \emph{native} Ackermann
   controller (server-side PID)};
\node[brg, below=\gap of ctrl] (sth)
  {\textbf{Per-sensor threads}\\[1pt]
   one publish thread each;\\optional dedicated clients};
\node[brg, below=\gap of sth] (tth)
  {\textbf{Telemetry thread}\\[1pt]
   10\,Hz on its own thread --- each\\tick costs blocking CARLA RPCs};
\node[brg, below=\gap of tth] (clk)
  {\textbf{\ttfamily sensor\_clock}\\[1pt]
   every stamp drawn from the\\simulation epoch, not wall time};

\begin{scope}[on background layer]
  \node[draw, dashed, rounded corners=2pt, inner sep=3.6pt, fill=black!1,
        fit=(ctrl)(sth)(tth)(clk)] (bridge) {};
\end{scope}
\node[colhead, above=1.4mm of bridge.north]
  {\textsc{Carla\_ASU\_Bridge}\\{\scriptsize\mdseries ROS 2 lifecycle node (C\texttt{++})}};

% ============================ 3. ROS 2 topics ============================
\node[top, anchor=north west] (tsub) at (11.1,0)
  {/drive\\control/tire\_friction\\control/drag\_coefficient};
\node[note, below=1.2mm of tsub] (nsub)
  {\emph{subscribed} --- commands and\\runtime physics overrides};

\node[top, below=2.6mm of nsub] (tstate)
  {/odom\\feedback/steering\_angle\\feedback/imu\\feedback/speed};
\node[note, below=1.2mm of tstate] (nstate)
  {\emph{published} --- the signals the\\identification actually consumes};

\node[gt, below=2.6mm of nstate] (tgt)
  {feedback/tire\_forces};
\node[note, below=1.2mm of tgt] (ngt)
  {\textbf{ground truth} --- per-wheel slip\\
   angle, normal load and lateral\\
   force, straight from PhysX; no\\
   analytic tire model in between};

\begin{scope}[on background layer]
  \node[draw, dashed, rounded corners=2pt, inner sep=3.6pt, fill=black!1,
        fit=(tsub)(nsub)(tstate)(nstate)(tgt)(ngt)] (topics) {};
\end{scope}
\node[colhead, above=1.4mm of topics.north]
  {ROS 2 topic layer\\{\scriptsize\mdseries all stamped from \texttt{/clock}}};

% =========================== 4. Autonomy stack ===========================
\node[cons, anchor=north west] (pp) at (16.6,0)
  {\textbf{\ttfamily pure\_pursuit}\\[1pt]
   model-free bootstrap; needs\\only the wheelbase};
\node[cons, below=\gap of pp] (sid)
  {\textbf{\ttfamily On-Track-SysID}\\[1pt]
   residual NN $\rightarrow$ virtual sweep\\$\rightarrow$ Pacejka fit};
\node[cons, below=11mm of sid] (mpc)
  {\textbf{\ttfamily mpc\_path\_tracking}\\[1pt]
   \texttt{acados} NMPC; holds the\\admissibility gates};
\node[cons, below=\gap of mpc] (mgr)
  {\textbf{\ttfamily adaptive\_}\\\textbf{\ttfamily controller\_manager}\\[1pt]
   PP\,$\rightarrow$\,MPC arbitration};
\node[cons, below=\gap of mgr] (bch)
  {\textbf{Benchmark nodes}\\[1pt]
   1/5/10-step prediction metrics,\\$F_y$ against ground truth};

\begin{scope}[on background layer]
  \node[draw, dashed, rounded corners=2pt, inner sep=3.6pt, fill=black!2,
        fit=(pp)(sid)(mpc)(mgr)(bch)] (stack) {};
\end{scope}
\node[colhead, above=1.4mm of stack.north]
  {Autonomy stack\\{\scriptsize\mdseries this workspace}};

% ================================ arrows =================================
% CARLA <-> bridge, over the C++ client library.
% The two column groups have different heights, so carla.east and bridge.west
% do NOT share a y. Pin both endpoints to one row (the "Sensor actors" row)
% with the |- / -| intersection syntax so the arrow is exactly horizontal.
\draw[flow, <->, line width=0.7pt]
      (carla.east |- sens) -- (bridge.west |- sens)
      node[lbl, above, pos=0.5, yshift=0.3mm]{\ttfamily libcarla\_client}
      node[lbl, below, pos=0.5, yshift=-0.3mm]{RPC + streams};

% CONVENTION: every inter-column arrow begins and ends on a dashed GROUP
% boundary, never on an inner box -- an arrow that pierces the dashed frame
% reads as if it bypassed the column. The row it belongs to is carried by the
% y coordinate (|- <box>), so the arrow still lines up with its own topic.

% bridge -> published topics
\draw[flow] (bridge.east |- tstate) -- (topics.west |- tstate)
      node[lbl, above, pos=0.5]{publish};
\draw[gtflow] (bridge.east |- tgt) -- (topics.west |- tgt);

% commands travel right to left, back into the bridge
\draw[flow] (topics.west |- tsub) -- (bridge.east |- tsub)
      node[lbl, above, pos=0.5]{apply};

% topics -> autonomy stack, and the /drive return path
\draw[flow] (topics.east |- tstate) -- (stack.west |- tstate)
      node[lbl, above, pos=0.5]{state feedback};
\draw[flow] (stack.west |- tsub) -- (topics.east |- tsub)
      node[lbl, above, pos=0.5]{\ttfamily /drive};

% ground truth reaches ONLY the scoring nodes: out of the topic column at the
% tire_forces row, down the gutter, into the stack at the benchmark row.
\coordinate (gtout)  at (topics.east |- tgt);
\coordinate (gtbend) at ($(gtout)+(0.6,0)$);
\draw[gtflow] (gtout) -- (gtbend) |- (stack.west |- bch);
\node[font=\tiny, orange!80!black, rotate=90, anchor=south, inner sep=1.4pt]
      at ($(gtbend)!0.5!(gtbend |- bch)$) {validation only};

% the gated model interface, inside the stack
\draw[gateflow] (sid.south) -- (mpc.north)
      node[lbl, pos=0.5, align=center]
      {\ttfamily mpc/update\_params\\\rmfamily\emph{gated}};

% ================================= legend ================================
% Sits under columns 1--2, which are the shorter pair.
\node[anchor=north west, font=\tiny, align=left, draw, rounded corners=1.5pt,
      inner sep=4pt, fill=white, text width=88mm]
      at ($(carla.south west)+(0,-0.5)$)
  {\tikz[baseline=-0.5ex]\draw[gtflow](0,0)--(0.5,0);\ \textbf{ground-truth path.}
   Used only to \emph{score} the identifier; never inside the identification loop.\\[1.5pt]
   \tikz[baseline=-0.5ex]\draw[gateflow](0,0)--(0.5,0);\ \textbf{gated interface.}
   An identified tire set is accepted only if its grip ceiling, understeer sign and
   critical speed clear the controller's requirements (Sec.~\ref{subsec:gates}).\\[1.5pt]
   \tikz[baseline=-0.5ex]\draw[flow](0,0)--(0.5,0);\ ROS 2 topic traffic.};
\end{tikzpicture}
}
\caption{Architecture of the \bridge{} ROS~2 interface. The bridge links
against CARLA's C\texttt{++} client, drives the ego vehicle through CARLA's
native Ackermann controller, and republishes state, sensors and per-wheel
telemetry on ROS~2 topics stamped from the simulation clock. The
\texttt{feedback/\allowbreak tire\_forces} path (highlighted) carries the
simulator's own per-wheel slip angle, normal load and lateral force, and is the
ground truth against which the identifier is scored.}
\label{fig:bridge}
\end{figure*}

\subsection{\bridge{}: interfaces and telemetry}
\label{subsec:bridge}

\bridge{}~\cite{carlaasubridge} is a C\texttt{++} ROS~2 lifecycle node linked
directly against CARLA's client library (Fig.~\ref{fig:bridge}). It spawns and
configures the ego vehicle, drives its inputs, and republishes simulator state
as standard ROS~2 topics; all identification traffic crosses this interface.

\subsubsection{Command path}
The bridge runs in \texttt{ackermann\_drive} mode, subscribing to a single
\texttt{ack\-er\-mann\_msgs/AckermannDriveStamped} on \texttt{/drive} and
forwarding speed and steering angle to CARLA's native Ackermann controller,
whose server-side cascaded PID loops are configured once at startup. No
client-side steering PID is present---the commanded road-wheel angle passes
through in radians---so the achieved/\allowbreak commanded ratio of
Table~\ref{tab:vehicle} is a property of the vehicle and its steering curve
rather than of a tuning choice in the bridge.

\subsubsection{Feedback path}
Table~\ref{tab:topics} lists the interface consumed by the identifier and its
validation. Three properties are specific to this architecture.

\begin{table}[t]
\caption{ROS~2 interface used by the identification and validation pipeline.
Namespace is \texttt{/sim} unless marked absolute.}
\label{tab:topics}
\centering
\scriptsize
\setlength{\tabcolsep}{3pt}
\begin{tabular}{@{}p{2.4cm}p{2.75cm}p{2.5cm}@{}}
\toprule
Topic & Type & Role\\
\midrule
\texttt{/odom} & \texttt{nav\_msgs/Odometry} & State $v_{x},v_{y},\omega$; twist is body-frame\\
\texttt{feedback/\allowbreak steering\_angle} & \texttt{std\_msgs/Float32} (deg) & Achieved road-wheel angle $\delta$\\
\texttt{feedback/imu} & \texttt{sensor\_msgs/Imu} & $a_{x},a_{y}$ for the friction warm-start\\
\texttt{feedback/\allowbreak tire\_forces} & \texttt{sim\_manager\_\allowbreak msgs/\allowbreak TireForces} & Per-wheel $\alpha$, $F_{z}$, $F_{y}$, $\mu$, $\omega_{w}$, $\kappa$: \textbf{ground truth}\\
\texttt{feedback/\allowbreak vehicle\_physics} & \texttt{std\_msgs/String} (JSON) & Live wheel parameters of \eqref{eq:physx}\\
\texttt{feedback/speed} & \texttt{std\_msgs/Float32} & Forward speed echo\\
\texttt{feedback/motors} & \texttt{std\_msgs/String} (JSON) & Per-wheel speed/torque/brake\\
\texttt{/drive} & \texttt{ackermann\_msgs/\allowbreak Ackermann\allowbreak Drive\allowbreak Stamped} & Control command in\\
\texttt{control/\allowbreak tire\_friction} & \texttt{std\_msgs/Float32} & Runtime grip override\\
\texttt{/clock} \emph{(abs.)} & \texttt{rosgraph\_msgs/\allowbreak Clock} & Simulation time base\\
\bottomrule
\end{tabular}
\end{table}

\paragraph{Per-wheel ground truth}
\texttt{feedback/\allowbreak tire\_forces} carries CARLA's own
\texttt{Wheel\-TelemetryData}: per-wheel lateral slip angle, normal load,
lateral force, effective friction coefficient, wheel speed, longitudinal force
and wheel torque, in wheel order [FL, FR, RL, RR]. No analytic tire model
exists inside the bridge, so nothing on this topic can disagree with the
physics the vehicle is driven by. This topic is what allows an on-track
identifier to be scored against the forces the plant generated, and the
channels it validates against are the lateral ones. Four properties of the
channel were established by measurement and bound what may be read off it.

$\bullet$~\emph{Lateral force is a contact-patch force.} In steady cornering
$\sum F_{y}$ tracks $m a_{y}$ with correlation \numrange{0.95}{0.97} and slope
\numrange{0.92}{1.08}, and $-F_{y}$ against $\alpha F_{z}$ gives
\SI{11.3}{\per\radian} (corr.\ \num{0.93}), a peak near \SI{5.3}{\degree} of
slip.

$\bullet$~\emph{Longitudinal force is not.} It behaves as a friction-capacity
report: over 8268 wheel samples $|F_{x}|$ sits above $0.80\,\mu F_{z}$ half the
time and exactly at $\mu F_{z}$ \SI{15}{\percent} of the time, while
$\mathrm{corr}(\sum F_{x},\,m a_{x})=\num{0.04}$; cutting engine torque by
\SI{89}{\percent} left it pinned at the same limit. Wheel torque is the
identity $-F_{x}r_{w}$ and carries nothing further. Neither channel is used
here, and the longitudinal capacity is taken as $\mu F_{z}$ instead.

$\bullet$~\emph{Normal load carries a known offset.} $\sum F_{z}$ runs
\SIrange{13}{15}{\percent} above $mg$ for the reasons given with
Table~\ref{tab:vehicle}, and its load transfer is only $\approx\SI{15}{\percent}$
of the rigid-body expectation while the configured centre of mass sits at
ground level---which is why the warm start of Section~\ref{subsec:friction}
builds $F_{z}$ from the measured static split and an explicit transfer term
rather than reading this field.

$\bullet$~\emph{Standstill is gated.} Below \SI{0.5}{\meter\per\second} the
slip, force and torque fields publish zero, because PhysX sleeps a parked
vehicle and then repeats the last telemetry frame byte for byte, and because
$\alpha$ is an $\arctan$ of two vanishing velocities: at
\SI{0.01}{\meter\per\second} it read \SIrange{-2.0}{2.6}{\degree} with
$\pm\SI{334}{\newton}$ of lateral force on wheels already at
$\omega_{w}=0$. The identification and benchmarking gates of
Sections~\ref{subsec:data} and~\ref{subsec:validation} sit well above this
speed for the same reason.

\paragraph{Simulation-time stamping}
Every feedback message is stamped from the same simulation epoch as
\texttt{/clock}. In synchronous mode CARLA ticks as fast as the server manages
(observed near $1.9\times$ real time), so wall-clock stamping would drift the
identification dataset against itself without bound.

\paragraph{Odometry and frame conventions}
Odometry is computed from the actor transform and velocity at
\SI{50}{\hertz} with \texttt{follow\_server\_rate} enabled and an optional
additive Gaussian noise model (disabled here). Publisher and subscribers all
use \texttt{BEST\_EFFORT}, because a reliability mismatch on this interface is
silent: DDS requires requested $\le$ offered, so a \texttt{RELIABLE} request
against a \texttt{BEST\_EFFORT} publisher yields one warning and then no data.
Because \eqref{eq:slip} is meaningless in the wrong frame, the body-frame
assumption on \texttt{/odom}'s twist was verified rather than assumed:
correlation \num{0.93} with \SI{0.09}{\meter\per\second} RMS error against the
reconstructed body-frame lateral velocity, versus
\SI{16.7}{\meter\per\second} RMS under a world-frame interpretation.

Figure~\ref{fig:frames} states the conventions the rest of the paper assumes.
The simulator and the identification stack do not share a handedness: CARLA
inherits Unreal's left-handed world (\mbox{$X$ east}, $Y$ \emph{south}, $Z$ up,
yaw positive clockwise), while every ROS~2 topic follows
REP-103~\cite{rep103}---right-handed ENU for the world, forward-left-up for the
body---and the identification equations of Section~\ref{sec:method} are written
in the body frame with $\alpha$ and $F_{y}$ positive to the left, the SAE-style
axis convention~\cite{saej670,pacejka2012book}. The bridge is where the two
meet, and it is a per-channel conversion rather than a single transform:
position and yaw flip sign, twist is additionally rotated into the body frame
by $q^{-1}$, the IMU's $a_{y}$ flips, the steering \emph{command} flips on the
way out, and the tire telemetry's slip angle does \emph{not}, because the frame
flip and CARLA's own sign convention for \texttt{lat\_slip} cancel. That last
one is the trap: a sign error there is invisible in the states and inverts the
identified curve, so it was checked against a slip angle rebuilt from the
chassis state in a steady \SI{4}{\meter\per\second} turn at
\SI{0.20}{\radian}---equal magnitude to within \SI{2}{\percent} front and
\SI{15}{\percent} rear.

\begin{figure*}[t]
\centering
\resizebox{\textwidth}{!}{% =============================================================================
% Fig. -- Frame and sign conventions across the CARLA/ROS 2 boundary.
%
% Vector TikZ, drawn at the paper's own font sizes. Included by
% sections/platform.tex inside a figure* (full text width).
%
% Three panels, left to right:
%   (a) the two world frames and the map between them (CARLA/Unreal is
%       left-handed with Y south and yaw clockwise; ROS is ENU/REP-103).
%   (b) the vehicle body frame in which every identification quantity is
%       defined: FLU, positive steer to the left, positive lateral force to
%       the left, slip angles as in (2).
%   (c) the per-channel conversion the bridge applies, which is where a sign
%       error would enter the identification silently.
%
% Layout note: every panel's bounding box is fitted to two explicit corner
% coordinates rather than to its content, so the three panels keep the same
% height and nothing collides once the whole picture is resized to \textwidth.
% =============================================================================

\begin{tikzpicture}[
  font=\scriptsize,
  >={Stealth[length=1.7mm,width=1.5mm]},
  ax/.style      = {->, line width=0.5pt, black!75},
  car/.style     = {line width=0.7pt, black!75},
  wheel/.style   = {line width=1.5pt, black!80, line cap=round},
  frc/.style     = {->, line width=0.7pt, green!45!black},
  vel/.style     = {->, line width=0.6pt, blue!55!black},
  rot/.style     = {->, line width=0.5pt, black!60},
  lbl/.style     = {font=\tiny, inner sep=1.0pt},
  note/.style    = {font=\tiny, align=left, inner sep=1.6pt},
  hdr/.style     = {font=\scriptsize\bfseries, align=center},
  conv/.style    = {->, line width=0.6pt, black!65, dashed}
]

% ================== (a) vehicle body frame and sign rules ==================
% Single-track sketch: x forward (right), y left (up), omega CCW positive.
\begin{scope}[shift={(93mm,2mm)}]
  \coordinate (cg)  at (0,0);
  \coordinate (fr)  at (15mm,0);
  \coordinate (rr)  at (-13mm,0);

  \draw[car] (rr) -- (fr);
  \draw[wheel] ($(rr)+(-3mm,0)$) -- ($(rr)+(3mm,0)$);
  \draw[wheel, rotate around={17:(fr)}] ($(fr)+(-3mm,0)$) -- ($(fr)+(3mm,0)$);
  \fill (cg) circle (0.6mm);

  % body axes at the centre of gravity
  \draw[ax] (cg) -- ++(8mm,0) node[lbl,below] {$x$ fwd};
  \draw[ax] (cg) -- ++(0,7mm) node[lbl,left]  {$y$ left};
  \draw[rot] (-2.8mm,-2.8mm) arc (225:135:4.0mm);
  \node[lbl, anchor=north east] at (-3.2mm,-2.4mm) {$\omega>0$};

  % steer angle at the front wheel
  \draw[black!45, line width=0.35pt] (fr) -- ++(6mm,0);
  \draw[rot] ($(fr)+(5mm,0)$) arc (0:17:5mm);
  \node[lbl, anchor=west] at ($(fr)+(5.4mm,2.2mm)$) {$\delta>0$};

  % velocity and sideslip at the CG
  \draw[vel] (cg) -- ++(10mm,4.5mm) node[lbl,above] {$v$};
  \draw[rot] (5.5mm,0) arc (0:24.2:5.5mm);
  \node[lbl, anchor=west] at (6.0mm,1.2mm) {$\beta$};

  % axle lateral forces, positive to the left
  \draw[frc] (fr) -- ++(0,7mm) node[lbl,above] {$F_{y,f}>0$};
  \draw[frc] (rr) -- ++(0,7mm) node[lbl,above] {$F_{y,r}>0$};
\end{scope}

\node[note, anchor=north west, text width=62mm] at (63mm,-11mm)
  {$\alpha_{f}=\delta-\arctan\dfrac{v_{y}+l_{f}\omega}{v_{x}}$,\qquad
   $\alpha_{r}=-\arctan\dfrac{v_{y}-l_{r}\omega}{v_{x}}$\\[2pt]
   Positive $\alpha$ produces positive (leftward) $F_{y}$. Wheel order in
   every per-wheel array: [FL, FR, RL, RR].};

\coordinate (pb1) at (61mm,-27mm);
\coordinate (pb2) at (127mm,15mm);
\begin{scope}[on background layer]
  \node[draw, dashed, rounded corners=2pt, inner sep=0pt, fill=black!2,
        fit=(pb1)(pb2)] (panelb) {};
\end{scope}
\node[hdr, above=1.0mm of panelb.north] {(a) Vehicle body frame (FLU)};

% ====================== (b) per-channel conversions =========================
\node[note, anchor=north west] at (133mm,13mm)
  {\renewcommand{\arraystretch}{1.2}%
   \begin{tabular}{@{}p{17mm}p{40mm}@{}}
   \textbf{Channel} & \textbf{Bridge conversion}\\
   \hline\\[-6pt]
   \texttt{/odom} pose & $y\mapsto-y$; $\psi,\theta\mapsto-\psi,-\theta$\\
   \texttt{/odom} twist & world $\to$ body by $q^{-1}$; $v_{y}$,
   $\omega_{y}$, $\omega_{z}$ negated\\
   \texttt{feedback/imu} & $a_{y}\mapsto-a_{y}$\\
   \texttt{tire\_forces} $\alpha$ & $\deg\to\mathrm{rad}$ only: the frame flip
   and CARLA's own sign convention cancel\\
   \texttt{tire\_forces} $F_{y}$ & $F_{y}\mapsto-F_{y}$, into $+y$ left\\
   \texttt{/drive} $\delta$ & negated on the way out to CARLA\\
   \end{tabular}};

\coordinate (pc1) at (131mm,-27mm);
\coordinate (pc2) at (196mm,15mm);
\begin{scope}[on background layer]
  \node[draw, dashed, rounded corners=2pt, inner sep=0pt, fill=black!2,
        fit=(pc1)(pc2)] (panelc) {};
\end{scope}
\node[hdr, above=1.0mm of panelc.north] {(b) Conversions applied by \bridge{}};

\end{tikzpicture}
}
\caption{Frame and sign conventions across the simulator--ROS~2 boundary.
(a) CARLA inherits Unreal's left-handed world frame with $Y$ pointing south and
yaw positive clockwise; ROS~2 uses right-handed ENU~\cite{rep103}, so the
bridge negates $y$, the yaw and pitch angles, and the corresponding angular
rates. (b) All identification quantities are defined in the body frame:
$x$ forward, $y$ left, $\omega$ and $\delta$ positive counter-clockwise, and
slip angle and lateral force positive to the left, so that a positive $\alpha$
produces a positive $F_{y}$~\cite{saej670}. (c) The conversion applied per
channel. The slip angle of the tire telemetry is the one exception that needs
no sign change, because the frame flip and CARLA's own convention cancel.}
\label{fig:frames}
\end{figure*}

\subsubsection{Extensions made for identification}
Six changes support this study and are listed for reproducibility: the
existing but unsubscribed \texttt{feedback/\allowbreak steering\_angle} is now
consumed by the identifier (with configurable units, scale and staleness
timeout, falling back to the \texttt{/drive} setpoint); runtime overrides
\texttt{control/\allowbreak tire\_friction} and
\texttt{control/\allowbreak drag\_coefficient} write straight to the live
\texttt{VehiclePhysicsControl}, so surface friction can change within a run
without respawning; the torque curve, zero-throttle damping and wheel
\texttt{tire\_friction} were tuned against a low-speed hold test (a wheel
friction above $\approx$2.0 on a steered wheel scrubs hard enough at creep
speed to stall the drivetrain intermittently), and these values define the
plant being identified; telemetry runs on a dedicated \SI{10}{\hertz}
thread off the control loop, since each tick costs several blocking CARLA RPCs;
\texttt{feedback/\allowbreak tire\_forces} was extended with the effective
friction coefficient, the wheel speed and a kinematically recomputed slip ratio
$\kappa=(\omega_{w}r_{w}-v)/v$, the last because CARLA's own
\texttt{long\_slip} field is inconsistent with the wheel speed and radius the
same message reports (\num{+0.177}, i.e.\ hard braking, on a wheel turning
\SI{8.99}{\radian\per\second} at a steady \SI{2.33}{\meter\per\second} cruise,
whose kinematic ratio is \num{-0.034}); and the live wheel parameters of
\eqref{eq:physx} are published on
\texttt{feedback/\allowbreak vehicle\_physics}, so the reference curve of
Table~\ref{tab:ref} cannot go stale when the vehicle configuration changes.

\subsection{Data acquisition and preprocessing}
\label{subsec:data}

The vehicle is driven around the reference raceline by a model-free
pure-pursuit controller, which requires no tire knowledge---only the
wheelbase---and is therefore a valid bootstrap, as in the baseline
\cite{ontrack2024,forzaeth2024}. Odometry, the achieved steering angle and the
drive command are sampled into a ring buffer at \SI{50}{\hertz}
($T_{s}=\SI{0.02}{\second}$) for \SI{30}{\second}, gated on
$v_{x}>\SI{1}{\meter\per\second}$. Preprocessing follows the baseline: a
zero-phase low-pass filter, then symmetry augmentation by mirroring
$(v_{y},\omega,\delta)\rightarrow(-v_{y},-\omega,-\delta)$ at unchanged
$v_{x}$, which doubles the data and removes the cornering bias of a circuit.
Re-identification retriggers every \SI{30}{\second} by default.

\subsection{How the architecture closes the loop}

The four stages of the identification pipeline map onto the architecture as
follows. \emph{Acquisition} uses \texttt{/odom},
\texttt{feedback/\allowbreak steering\_angle} and \texttt{/drive}, all stamped
on the simulation clock, so that the sampling interval used by the one-step
model is the interval the physics actually advanced. \emph{Preprocessing} and
\emph{fitting} run inside the identification node against that buffer.
\emph{Validation} draws on \texttt{feedback/\allowbreak tire\_forces} and
\texttt{feedback/imu}, which never enter the identification loop, so the
estimator is scored against a channel it did not observe. \emph{Deployment}
pushes the fitted coefficients over a ROS~2 service to the controller behind
the gates of Section~\ref{sec:integration}. Only the last stage leaves the
architecture, which is what makes the negative result of
Section~\ref{subsec:negative} a statement about the trajectory rather than
about the interface.
